% Section II: Foundations recap. Short — reproduces only what is load-bearing
% for the deployment and the empirical analysis. Full proofs and derivations
% are in the foundation paper \cite{lcp2025}. Figure captions B1, B2, B4
% drafted in References/Artifact_discussions.md §II.

This section recaps the load-bearing components of~\cite{lcp2025}: the support-normal equivalence theorem (Theorem~\ref{thm:equivalence}), the polyhedral $L_1$ and asymptotic $L_2$ specialisations (Propositions~\ref{prop:l1-polyhedral}--\ref{prop:l2-asymptotic}), the calibration algorithms (Algorithms~\ref{alg:alg1a}--\ref{alg:alg1b}), and the offline/online operational workflow. Full proofs, worked numerical examples, and the failure-mode catalogue are deferred to~\cite{lcp2025} and cited in place.

\subsection{Problem Statement}
\label{sec:foundations:problem}

We consider the discrete-time dynamical system $x_{k+1} = f(x_k, u_k)$ over horizon $N$, with state $x_k \in \mathbb{R}^n$ and control $u_k \in \mathbb{R}^m$. The trajectory is $z := (x_{0:N}, u_{0:N-1}) \in \mathbb{R}^d$ where $d := n(N+1) + mN$. The performance objective is
\begin{equation}
J(z) := \sum_{k=0}^{N-1} \ell(x_k, u_k) + \ell_f(x_N),
\label{eq:perf-objective}
\end{equation}
with stage cost $\ell$ and terminal cost $\ell_f$. Rule-based constraints are partitioned into $L$ priority levels $\mathcal{C}_i = \{z : g_{i,j}(z) \leq 0, \ j = 1, \ldots, m_i\}$, $i = 1, \ldots, L$, ordered $\mathcal{C}_1 \succ \mathcal{C}_2 \succ \cdots \succ \mathcal{C}_L$. The violation functional at level $i$ is
\begin{equation}
V_i(z) := \sum_{j=1}^{m_i} \bigl[g_{i,j}(z)\bigr]_{+}^{p}, \qquad p \in \{1, 2\},
\label{eq:violation-functional}
\end{equation}
with $p = 1$ the $L_1$ (hinge) penalty and $p = 2$ the $L_2$ (squared-hinge) penalty. The standing assumptions of~\cite{lcp2025} --- A1 existence of a lex cascade solution, A2 lex uniqueness up to performance-equivalent tie-breaks, and A3 convexity of $J$ and each $V_i$ --- are assumed throughout. Proposition~\ref{prop:l1-polyhedral} below additionally invokes a cascade Linear Independence Constraint Qualification (LICQ) regularity condition; failure of the cascade LICQ is one of the two structural failure modes catalogued in~\cite[\S8]{lcp2025} and audited empirically in Section~\ref{sec:results}~(\S\ref{sec:results:framework}).

The lexicographic cascade defines, with $\mathcal{F}_0 := \mathcal{Z}$ the dynamics-feasible set,
\begin{equation}
V_i^{\star} := \min_{z \in \mathcal{F}_{i-1}} V_i(z), \quad \mathcal{F}_i := \{z \in \mathcal{F}_{i-1} : V_i(z) = V_i^{\star}\},
\label{eq:cascade}
\end{equation}
for $i = 1, \ldots, L$, and minimises $J$ over $\mathcal{F}_L$ at stage $L+1$. We call the resulting trajectory $z_{\mathrm{lex}}^{\star}$ and its image $p^{\star} = (V_1^{\star}, \ldots, V_L^{\star}, J(z_{\mathrm{lex}}^{\star}))$. The weighted-sum scalarisation produces $z_{\mathrm{ws}}^{\star}(w) := \arg\min_z\, [J(z) + \sum_i w_i V_i(z)]$ in a single solve, parameterised by the weight vector $w \in \mathbb{R}_{++}^{L}$.

\subsection{Support-Normal Equivalence}
\label{sec:foundations:thm1}

Let $\overline{\mathcal{P}}$ denote the upper image of the achievement map $\Phi(z) = (V_1(z), \ldots, V_L(z), J(z))$. Let $\widehat{\Omega}(p^{\star}) \subset \mathbb{R}^{L+1}$ denote the outward normal cone of $\overline{\mathcal{P}}$ at $p^{\star}$, and $\Omega(p^{\star}) \subset \mathbb{R}_{++}^{L}$ its unit-performance slice. The main theorem of~\cite{lcp2025} states:

\begin{theorem}[Support-normal equivalence~{\cite[Thm.~4.1]{lcp2025}}]
\label{thm:equivalence}
The lex optimum $z_{\mathrm{lex}}^{\star}$ is contained in $\arg\min_z\, [J(z) + \sum_i w_i V_i(z)]$ if and only if $w \in \Omega(p^{\star})$. Equality with the selected weighted-sum solution holds under weighted-sum uniqueness or under an explicit lex-preserving tie-break.
\end{theorem}

The proof in~\cite[\S4.2]{lcp2025} proceeds by noting that the lex point $p^\star$ is an extreme point of the upper image $\overline{\mathcal{P}}$ (Lemma~3.2 of~\cite{lcp2025}); the weighted-sum solve at $w$ selects the boundary point of $\overline{\mathcal{P}}$ exposed by the affine functional $(w, 1)$, which coincides with $p^\star$ exactly when $(w, 1)$ supports $\overline{\mathcal{P}}$ at $p^\star$. The condition is both necessary and sufficient: substituting the calibrated weighted-sum solve for the cascade is operationally licensed only when the weight vector lies in $\Omega(p^\star)$, which is precisely the calibration problem Algorithm~\ref{alg:alg1a} solves. The achievement map $\Phi$, the upper image $\overline{\mathcal{P}}$, and the lex image $p^\star$ are reified in our reference implementation as the \texttt{lcp.AchievementImage} dataclass and the \texttt{lcp.lex\_image\_from\_cascade} and \texttt{lcp.verify\_lex\_extreme\_point} helpers; the latter provides a numerical extreme-point check on a finite probe set.

\subsection{Polyhedral $L_1$ Regime}
\label{sec:foundations:l1}

For $L_1$ violation functionals (\eqref{eq:violation-functional} with $p = 1$), the upper image is polyhedral and the equivalence region admits explicit characterisation:

\begin{proposition}[$L_1$ polyhedral structure~{\cite[Thm.~5.1]{lcp2025}}]
\label{prop:l1-polyhedral}
Under the standing assumptions and Setting~(P) of~\cite{lcp2025}, $\overline{\mathcal{P}}$ is a full-dimensional polyhedron; $p^{\star}$ is a vertex; and $\Omega(p^{\star})$ is a non-empty, full-dimensional polyhedral region in $\mathbb{R}_{++}^{L}$.
\end{proposition}

The full-dimensionality of $\Omega(p^{\star})$ is the geometric content of the empirical observation that ``moderate weights work''. \emph{Remark.} Cascade Linear Independence Constraint Qualification (LICQ) is \emph{not} a hypothesis of Proposition~\ref{prop:l1-polyhedral}; the geometric statements follow from the polyhedrality of $\overline{\mathcal{P}}$ and the extreme-point property of $p^{\star}$ alone (cf.~\cite[Remark~5.1]{lcp2025}). Cascade LICQ is required at the algorithmic step where the explicit half-space description of $\Omega(p^{\star})$ is computed by Fourier--Motzkin elimination of the lex KKT multipliers in Algorithm~\ref{alg:alg1a} (Step~3); LICQ failure makes $\Omega(p^{\star})$ lower-dimensional but not empty (audited empirically in Section~\ref{sec:results:framework}).

\subsection{$L_2$ Asymptotic Regime}
\label{sec:foundations:l2}

For $L_2$ violation functionals ($p = 2$), the squared-hinge penalty $[g]_+^2$ has $\nabla [g]_+^2 = 2[g]_+ \nabla g = 0$ at every boundary point ($g = 0$). Consequently no finite weight produces a first-order force that resolves the lex priority at the constraint surface, and exact equivalence in the sense of Theorem~\ref{thm:equivalence} is generically unavailable. The operationally relevant notion is asymptotic compliance:

\begin{proposition}[$L_2$ asymptotic rates~{\cite[Thm.~6.1]{lcp2025}}]
\label{prop:l2-asymptotic}
Suppose Setting~(Q) and the standing assumptions hold, and that the lex optimum $z_{\mathrm{lex}}^{\star}$ satisfies the four local regularity conditions: (a)~the active set $\mathcal{A}^{\star}$ is stable in a neighbourhood of $z_{\mathrm{lex}}^{\star}$ under small perturbations of weights; (b)~the gradients $\nabla g_{i, j}(z_{\mathrm{lex}}^{\star})$ at boundary-binding constraints are non-zero and linearly independent; (c)~second-order sufficient conditions hold (the reduced Hessian of $J$ on the active manifold is positive definite); and (d)~the weighted-sum minimiser is unique in a neighbourhood of $z_{\mathrm{lex}}^{\star}$ for weights in the regime considered. Suppose additionally the performance-outward sign condition $\langle \nabla J(z_{\mathrm{lex}}^{\star}), \nabla g_{i, j}(z_{\mathrm{lex}}^{\star}) \rangle < 0$ at the boundary-binding constraint $g_{i, j}$. Then the raw constraint violation and the violation functional decay at the rates
\begin{equation}
    g_{i, j}(z_{\mathrm{ws}}^{\star}(w)) = O(1/w_i), \qquad V_i(z_{\mathrm{ws}}^{\star}(w)) = O(1/w_i^2)
    \label{eq:l2-asymp}
\end{equation}
as $w_i \to \infty$.
\end{proposition}

Equation~\eqref{eq:l2-asymp} is consistent with the $O(1/\sqrt{\epsilon_i})$ threshold scaling of compliance equivalence (Section~\ref{sec:foundations:compliance}): inverting $V_i(z_{\mathrm{ws}}^{\star}(w)) = O(1/w_i^2) \leq \epsilon_i$ gives $w_i = O(1/\sqrt{\epsilon_i})$ for the boundary-binding regime in Algorithm~\ref{alg:alg1b}.

\subsection{Compliance Equivalence}
\label{sec:foundations:compliance}

Where exact equivalence is unavailable, the operationally relevant notion is preservation of the binary satisfaction pattern. Let $b_{\epsilon}(z) \in \{0, 1\}^{L}$ denote the per-level compliance vector at tolerance vector $\boldsymbol{\epsilon} = (\epsilon_1, \ldots, \epsilon_L)$, with $b_{\epsilon, i}(z) = \mathbb{1}[V_i(z) \leq \epsilon_i]$. The compliance-equivalence region is $\Omega_b^{\epsilon}(p^{\star}) := \{w \in \mathbb{R}_{++}^{L} : b_{\epsilon}(z_{\mathrm{ws}}^{\star}(w)) = b_{\epsilon}(z_{\mathrm{lex}}^{\star})\}$. The exact ($L_1$) and tolerance-based ($L_2$) versions of compliance equivalence and the threshold scaling $W_i(\epsilon_i, w_{-i}) = O(1/\sqrt{\epsilon_i})$ for squared-violation tolerance are stated in~\cite[\S7]{lcp2025}; we use compliance as the runtime safety net of Section~\ref{sec:deployment:compliance}. The compliance vector $b_{\epsilon}(z)$ is the multi-level generalisation of Hasuo's instantly-checkable RSS safety condition $C$~\cite[Def.~II.1, Req.~II.2]{hasuo2022rss}: where RSS specifies a single boolean condition per scenario together with a hand-designed proper response that maintains it, our framework specifies one boolean per priority level whose conjunction is the runtime certificate of lex-equivalence (Theorem~\ref{thm:equivalence}), with the cascade fallback (Section~\ref{sec:deployment:dual-mode}) playing the role of Hasuo's proper response uniformly across all levels.

\subsection{Algorithm 1A: $L_1$ Exact-Equivalence Calibration}
\label{sec:foundations:alg1a}

\begin{figure}[t]
    \centering
    % Fig.2 (B1) — Algorithm 1A flow (L_1 exact-equivalence calibration).
% Vertical column, 6 nodes, all sized to 7.0 cm wide for column alignment.
% Single-column (IEEE \columnwidth).

% Shared TikZ styles for the block-diagram figures (Fig.2-Fig.6). Loaded by
% each individual fig*.tex via \input. Each style sets BOTH minimum width AND
% text width explicitly so that every box in a column has the same width
% regardless of content length — this prevents the row-misalignment that
% happens when math content is wider than the nominal minimum width.

\tikzset{
  % ===== SINGLE-COLUMN FLOW (Fig.2, Fig.3, Fig.4) =====
  % Width 7.0 cm fits IEEE single-column (\columnwidth ~ 8.4 cm).
  colblock/.style = {
    rectangle, draw=black!70, line width=0.35pt, rounded corners=2pt,
    text width=6.45cm, align=center, font=\sffamily\footnotesize,
    minimum width=6.9cm, minimum height=0.95cm,
    inner sep=5pt, outer sep=0pt, fill=white
  },
  colblockio/.style = {
    colblock,
    fill=black!3
  },
  coldecision/.style = {
    diamond, draw=black!70, line width=0.35pt, aspect=2.5,
    text width=4.1cm, align=center, font=\sffamily\footnotesize,
    inner sep=1.5pt, outer sep=0pt, fill=white
  },
  % ===== DOUBLE-COLUMN GRID (Fig.5) =====
  % Used in figure* environments, ~14 cm wide.
  gridcell/.style = {
    rectangle, draw=black!70, line width=0.35pt, rounded corners=1.5pt,
    text width=2.8cm, align=center, font=\sffamily\scriptsize,
    minimum width=3.0cm, minimum height=0.55cm,
    inner sep=2.5pt, outer sep=0pt, fill=white
  },
  gridcellstub/.style = {
    gridcell,
    draw=black!40, text=black!55
  },
  % ===== DOUBLE-COLUMN HORIZONTAL FLOW (Fig.6) =====
  hblock/.style = {
    rectangle, draw=black!70, line width=0.35pt, rounded corners=2pt,
    text width=2.5cm, align=center, font=\sffamily\scriptsize,
    minimum width=2.7cm, minimum height=1.0cm,
    inner sep=4pt, outer sep=0pt, fill=white
  },
  hblockio/.style = {
    hblock,
    fill=black!3
  },
  % ===== COMMON =====
  phasehdr/.style = {
    font=\sffamily\small\bfseries, text=black!78,
    align=center
  },
  boxhdr/.style = {
    font=\sffamily\bfseries\footnotesize,
    align=center,
    text=black
  },
  boxmeta/.style = {
    font=\sffamily\scriptsize,
    align=center,
    text=black!78
  },
  semioffline/.style = {fill=blue!5},
  semianalysis/.style = {fill=orange!8},
  seminovel/.style = {fill=green!10, draw=black!80, line width=0.5pt},
  semiruntime/.style = {fill=gray!7},
  semidecision/.style = {fill=orange!12},
  semiwarn/.style = {fill=red!8},
  flow/.style = {
    -Stealth, semithick, draw=black
  },
  feedback/.style = {
    -Stealth, semithick, draw=black!65, dashed
  },
  arrowlbl/.style = {
    font=\sffamily\scriptsize,
    fill=white, inner sep=1.5pt
  },
  % Translucent group background
  groupbg/.style = {
    rectangle, draw=black!30, line width=0.35pt,
    rounded corners=3pt, inner sep=6pt, fill=white
  }
}

\begin{tikzpicture}[node distance=3.5mm]
  \node[colblockio, semioffline] (input) {%
    \textbf{Priority-ordered convex program}\\[-1pt]
    \scriptsize feasible weight box $\mathcal{W} \subset \mathbb{R}_{++}^{L}$
  };
  \node[colblock, semianalysis, below=of input] (lex) {%
    \textbf{1. Lex reference solve}\\[-1pt]
    \scriptsize outputs $(z_{\mathrm{lex}}^{\star},\ p^{\star})$
  };
  \node[colblock, semianalysis, below=of lex] (kkt) {%
    \textbf{2. KKT lifting}\\[-1pt]
    \scriptsize form the lex KKT system and lift to $(w,\beta)$-space
  };
  \node[colblock, semianalysis, below=of kkt] (fm) {%
    \textbf{3. Equivalence-region projection}\\[-1pt]
    \scriptsize Fourier--Motzkin yields $\Omega(p^{\star})$ in weight space
  };
  \node[colblock, seminovel, below=of fm] (lp) {%
    \textbf{4. Weight calibration}\\[-1pt]
    \scriptsize Chebyshev-centre LP on $\Omega(p^{\star}) \cap \mathcal{W}$
  };
  \node[colblockio, semiruntime, below=of lp] (out) {%
    \textbf{Offline artefacts}\\[-1pt]
    \scriptsize calibrated weight $w^{\dagger}$ and reference compliance $b_{\epsilon}^{\star}$
  };

  \draw[flow] (input.south) -- (lex.north);
  \draw[flow] (lex.south)   -- (kkt.north);
  \draw[flow] (kkt.south)   -- (fm.north);
  \draw[flow] (fm.south)    -- (lp.north);
  \draw[flow] (lp.south)    -- (out.north);
\end{tikzpicture}

    \caption{Algorithm~\ref{alg:alg1a} as an offline calibration workflow for the $L_1$ regime. The lex reference solve provides the active-set signature, the KKT and projection steps recover the exact-equivalence region in weight space, and the Chebyshev-centre program selects a robust deployment weight $w^{\dagger}$ inside the admissible box. Cf.~\cite[Fig.~B1]{lcp2025}.}
    \label{fig:alg1a_flow}
\end{figure}

\begin{algorithm}[t]
\caption{$L_1$ exact-equivalence calibration~\cite[\S9.3]{lcp2025}}
\label{alg:alg1a}
\begin{algorithmic}[1]
\Require Convex priority-ordered program with $L_1$ violation functionals; weight box $\mathcal{W} = [\underline{w}_i, \overline{w}_i]^L$.
\State Solve the lex cascade~\eqref{eq:cascade} to obtain $z_{\mathrm{lex}}^{\star}$ and the active-constraint signature $p^{\star}$.
\State Form the lex Karush--Kuhn--Tucker (KKT) system at $z_{\mathrm{lex}}^{\star}$ and lift to the $(w, \beta)$-space.
\State Eliminate the per-level lex multipliers via Fourier--Motzkin to obtain the half-space description of $\Omega(p^{\star})$.
\State Solve the box-bounded Chebyshev-centre linear program on $\Omega(p^{\star}) \cap \mathcal{W}$ to obtain $w^{\dagger}$ and inscribed-ball radius $r^{\dagger}$.
\Ensure Calibrated weight vector $w^{\dagger}$; cache the reference compliance vector $b_{\epsilon}^{\star} := b_{\epsilon}(z_{\mathrm{lex}}^{\star})$.
\end{algorithmic}
\end{algorithm}

Algorithm~\ref{alg:alg1a} operates in the polyhedral $L_1$ regime (Proposition~\ref{prop:l1-polyhedral}) and returns a weight vector $w^{\dagger}$ at which the weighted-sum solve provably recovers the lex optimum. Worked Example~1 of~\cite[\S11.1]{lcp2025} reproduces analytically $w^{\dagger} = (5.5, 5.5)$ and $r^{\dagger} = 4.5$ on a two-priority instance; the algorithm is unit-tested in our codebase against this reference.

The framework also provides Algorithm~0 (the homogeneous-cone primary formulation of~\cite[\S9.1]{lcp2025}), an operator-supplied-box-free alternative to Algorithm~1A that solves a Chebyshev linear program on the simplex slice of the homogeneous cone $\widehat\Omega(p^\star)$ and projects back to unit-performance weights via $\bar w_i := w_i^\sharp / \beta^\sharp$. On Example~1 it returns $\bar w \approx (2, 2)$, also inside $\Omega(p^\star)$. The detailed pseudocode and complexity bounds for Algorithms~0, 1A, and 1B are in Appendix~\ref{app:algorithms}.

\subsection{Algorithm 1B: $L_2$ Tolerance-Compliance Calibration}
\label{sec:foundations:alg1b}

\begin{figure}[t]
    \centering
    % Fig.3 (B2) — Algorithm 1B flow (L_2 tolerance-compliance calibration).
% Vertical column with a decision diamond and a feedback loop routed around
% the right side of the column (no crossing of intermediate boxes).
% Single-column (IEEE \columnwidth ~ 8.4 cm).
%
% Width budget: column boxes are 7.0 cm wide (centred); feedback path runs
% along x = +4.0 cm to +4.6 cm to stay clear of the box's right edge
% (which extends to x = +3.5 cm).

% Shared TikZ styles for the block-diagram figures (Fig.2-Fig.6). Loaded by
% each individual fig*.tex via \input. Each style sets BOTH minimum width AND
% text width explicitly so that every box in a column has the same width
% regardless of content length — this prevents the row-misalignment that
% happens when math content is wider than the nominal minimum width.

\tikzset{
  % ===== SINGLE-COLUMN FLOW (Fig.2, Fig.3, Fig.4) =====
  % Width 7.0 cm fits IEEE single-column (\columnwidth ~ 8.4 cm).
  colblock/.style = {
    rectangle, draw=black!70, line width=0.35pt, rounded corners=2pt,
    text width=6.45cm, align=center, font=\sffamily\footnotesize,
    minimum width=6.9cm, minimum height=0.95cm,
    inner sep=5pt, outer sep=0pt, fill=white
  },
  colblockio/.style = {
    colblock,
    fill=black!3
  },
  coldecision/.style = {
    diamond, draw=black!70, line width=0.35pt, aspect=2.5,
    text width=4.1cm, align=center, font=\sffamily\footnotesize,
    inner sep=1.5pt, outer sep=0pt, fill=white
  },
  % ===== DOUBLE-COLUMN GRID (Fig.5) =====
  % Used in figure* environments, ~14 cm wide.
  gridcell/.style = {
    rectangle, draw=black!70, line width=0.35pt, rounded corners=1.5pt,
    text width=2.8cm, align=center, font=\sffamily\scriptsize,
    minimum width=3.0cm, minimum height=0.55cm,
    inner sep=2.5pt, outer sep=0pt, fill=white
  },
  gridcellstub/.style = {
    gridcell,
    draw=black!40, text=black!55
  },
  % ===== DOUBLE-COLUMN HORIZONTAL FLOW (Fig.6) =====
  hblock/.style = {
    rectangle, draw=black!70, line width=0.35pt, rounded corners=2pt,
    text width=2.5cm, align=center, font=\sffamily\scriptsize,
    minimum width=2.7cm, minimum height=1.0cm,
    inner sep=4pt, outer sep=0pt, fill=white
  },
  hblockio/.style = {
    hblock,
    fill=black!3
  },
  % ===== COMMON =====
  phasehdr/.style = {
    font=\sffamily\small\bfseries, text=black!78,
    align=center
  },
  boxhdr/.style = {
    font=\sffamily\bfseries\footnotesize,
    align=center,
    text=black
  },
  boxmeta/.style = {
    font=\sffamily\scriptsize,
    align=center,
    text=black!78
  },
  semioffline/.style = {fill=blue!5},
  semianalysis/.style = {fill=orange!8},
  seminovel/.style = {fill=green!10, draw=black!80, line width=0.5pt},
  semiruntime/.style = {fill=gray!7},
  semidecision/.style = {fill=orange!12},
  semiwarn/.style = {fill=red!8},
  flow/.style = {
    -Stealth, semithick, draw=black
  },
  feedback/.style = {
    -Stealth, semithick, draw=black!65, dashed
  },
  arrowlbl/.style = {
    font=\sffamily\scriptsize,
    fill=white, inner sep=1.5pt
  },
  % Translucent group background
  groupbg/.style = {
    rectangle, draw=black!30, line width=0.35pt,
    rounded corners=3pt, inner sep=6pt, fill=white
  }
}

\begin{tikzpicture}[node distance=3.5mm]
  \node[colblockio, semioffline] (input) {%
    \textbf{$L_2$ deployment target}\\[-1pt]
    \scriptsize tolerance vector $\boldsymbol{\epsilon}$ and feasible weight box $\mathcal{W}$
  };
  \node[colblock, semianalysis, below=of input] (lex) {%
    \textbf{1. Lex reference solve}\\[-1pt]
    \scriptsize outputs $(z_{\mathrm{lex}}^{\star},\ p^{\star})$
  };
  \node[colblock, semianalysis, below=of lex] (sens) {%
    \textbf{2. Sensitivity extraction}\\[-1pt]
    \scriptsize singular-perturbation coefficients $(c_i,\kappa_i)$ at $p^{\star}$
  };
  \node[colblock, semianalysis, below=of sens] (thresh) {%
    \textbf{3. Threshold synthesis}\\[-1pt]
    \scriptsize admissible lower bounds $W_i(\epsilon_i, w_{-i})$
  };
  \node[colblock, seminovel, below=of thresh] (lp) {%
    \textbf{4. Weight calibration}\\[-1pt]
    \scriptsize coupled Chebyshev-centre LP under the threshold constraints
  };
  \node[coldecision, semidecision, below=6mm of lp] (verify) {%
    Compliance preserved?\\[-1pt]
    $b_{\epsilon}(z_{\mathrm{ws}}^{\star}(w^{\dagger})) = b_{\epsilon}^{\star}$
  };
  \node[colblockio, semiruntime, below=6mm of verify] (out) {%
    \textbf{Offline artefacts}\\[-1pt]
    \scriptsize calibrated weight $w^{\dagger}$ and cached reference $b_{\epsilon}^{\star}$
  };

  \draw[flow] (input.south)  -- (lex.north);
  \draw[flow] (lex.south)    -- (sens.north);
  \draw[flow] (sens.south)   -- (thresh.north);
  \draw[flow] (thresh.south) -- (lp.north);
  \draw[flow] (lp.south)     -- (verify.north);
  \draw[flow] (verify.south) -- node[arrowlbl, right=2pt] {match} (out.north);

  % Feedback (mismatch) — route around the right side, between thresh and
  % verify, well clear of any box's right edge (boxes extend to x ≈ +3.5 cm).
  \coordinate (rEast)   at ([xshift=8mm] verify.east);
  \coordinate (rThresh) at ([xshift=8mm] thresh.east);
  \draw[feedback]
    (verify.east) -- (rEast)
    node[arrowlbl, above=1pt, midway] {mismatch: tighten tolerance}
    -- (rThresh) -- (thresh.east);
\end{tikzpicture}

    \caption{Algorithm~\ref{alg:alg1b} as an offline calibration workflow for the $L_2$ regime. Exact finite-weight equivalence is replaced by compliance preservation: the lex reference is analysed to obtain sensitivity coefficients, the threshold step synthesises admissible lower bounds $W_i(\epsilon_i, w_{-i})$, and the coupled calibration program selects $w^{\dagger}$. A failed compliance check sends the procedure back to tolerance refinement rather than silently deploying an invalid weight. Cf.~\cite[Fig.~B2]{lcp2025}.}
    \label{fig:alg1b_flow}
\end{figure}

\begin{algorithm}[t]
\caption{$L_2$ tolerance-compliance calibration~\cite[\S9.4]{lcp2025}}
\label{alg:alg1b}
\begin{algorithmic}[1]
\Require Convex priority-ordered program with $L_2$ violation functionals; per-level tolerance $\boldsymbol{\epsilon}$; weight box $\mathcal{W}$.
\State Solve the lex cascade~\eqref{eq:cascade} to obtain $z_{\mathrm{lex}}^{\star}$ and $p^{\star}$.
\State Extract the singular-perturbation sensitivity constants $\{c_i, \kappa_i\}_{i=1}^{L}$ at $p^{\star}$.
\State Form the pointwise threshold $W_i(\epsilon_i, w_{-i}) = (c_i + \kappa_i^{T} w_{-i}) / \sqrt{\epsilon_i}$.
\State Solve the coupled-linear Chebyshev-centre linear program in $(w, \beta)$-space on $\{w \in \mathcal{W} : w_i \geq W_i(\epsilon_i, w_{-i}),\ i = 1, \ldots, L\}$ to obtain $w^{\dagger}$.
\State Verify the resulting compliance vector $b_{\epsilon}(z_{\mathrm{ws}}^{\star}(w^{\dagger}))$ matches $b_{\epsilon}^{\star}$; on mismatch, refine $\boldsymbol{\epsilon}$ or escalate.
\Ensure Calibrated weight vector $w^{\dagger}$; cached reference compliance vector $b_{\epsilon}^{\star}$.
\end{algorithmic}
\end{algorithm}

Worked Example~2 of~\cite[\S11.2]{lcp2025} reproduces analytically the threshold $W_1(\epsilon_1 = 0.01, w_3 = 1) \approx 77.5$ on a three-priority discrete-time instance; our codebase unit-tests Algorithm~\ref{alg:alg1b} against this reference.

\subsection{Operational Workflow}
\label{sec:foundations:state-machine}

\begin{figure}[t]
    \centering
    % Fig.4 (B4) — Offline/online operational workflow.
% Two grouped phases with an explicit runtime decision: accept on compliance
% match, otherwise log and fall back to the lex cascade.
% Single-column (IEEE \columnwidth ~ 8.4 cm).

% Shared TikZ styles for the block-diagram figures (Fig.2-Fig.6). Loaded by
% each individual fig*.tex via \input. Each style sets BOTH minimum width AND
% text width explicitly so that every box in a column has the same width
% regardless of content length — this prevents the row-misalignment that
% happens when math content is wider than the nominal minimum width.

\tikzset{
  % ===== SINGLE-COLUMN FLOW (Fig.2, Fig.3, Fig.4) =====
  % Width 7.0 cm fits IEEE single-column (\columnwidth ~ 8.4 cm).
  colblock/.style = {
    rectangle, draw=black!70, line width=0.35pt, rounded corners=2pt,
    text width=6.45cm, align=center, font=\sffamily\footnotesize,
    minimum width=6.9cm, minimum height=0.95cm,
    inner sep=5pt, outer sep=0pt, fill=white
  },
  colblockio/.style = {
    colblock,
    fill=black!3
  },
  coldecision/.style = {
    diamond, draw=black!70, line width=0.35pt, aspect=2.5,
    text width=4.1cm, align=center, font=\sffamily\footnotesize,
    inner sep=1.5pt, outer sep=0pt, fill=white
  },
  % ===== DOUBLE-COLUMN GRID (Fig.5) =====
  % Used in figure* environments, ~14 cm wide.
  gridcell/.style = {
    rectangle, draw=black!70, line width=0.35pt, rounded corners=1.5pt,
    text width=2.8cm, align=center, font=\sffamily\scriptsize,
    minimum width=3.0cm, minimum height=0.55cm,
    inner sep=2.5pt, outer sep=0pt, fill=white
  },
  gridcellstub/.style = {
    gridcell,
    draw=black!40, text=black!55
  },
  % ===== DOUBLE-COLUMN HORIZONTAL FLOW (Fig.6) =====
  hblock/.style = {
    rectangle, draw=black!70, line width=0.35pt, rounded corners=2pt,
    text width=2.5cm, align=center, font=\sffamily\scriptsize,
    minimum width=2.7cm, minimum height=1.0cm,
    inner sep=4pt, outer sep=0pt, fill=white
  },
  hblockio/.style = {
    hblock,
    fill=black!3
  },
  % ===== COMMON =====
  phasehdr/.style = {
    font=\sffamily\small\bfseries, text=black!78,
    align=center
  },
  boxhdr/.style = {
    font=\sffamily\bfseries\footnotesize,
    align=center,
    text=black
  },
  boxmeta/.style = {
    font=\sffamily\scriptsize,
    align=center,
    text=black!78
  },
  semioffline/.style = {fill=blue!5},
  semianalysis/.style = {fill=orange!8},
  seminovel/.style = {fill=green!10, draw=black!80, line width=0.5pt},
  semiruntime/.style = {fill=gray!7},
  semidecision/.style = {fill=orange!12},
  semiwarn/.style = {fill=red!8},
  flow/.style = {
    -Stealth, semithick, draw=black
  },
  feedback/.style = {
    -Stealth, semithick, draw=black!65, dashed
  },
  arrowlbl/.style = {
    font=\sffamily\scriptsize,
    fill=white, inner sep=1.5pt
  },
  % Translucent group background
  groupbg/.style = {
    rectangle, draw=black!30, line width=0.35pt,
    rounded corners=3pt, inner sep=6pt, fill=white
  }
}

\begin{tikzpicture}[node distance=3.8mm]
  \node[phasehdr] (offlinehdr) {Offline calibration};
  \node[colblockio, semioffline, below=2mm of offlinehdr] (scenario) {%
    \textbf{Scenario-class exemplar}\\[-1pt]
    \scriptsize select a representative operating case
  };
  \node[colblock, semianalysis, below=of scenario] (cascade) {%
    \textbf{Lex reference solve}\\[-1pt]
    \scriptsize obtain $(z_{\mathrm{lex}}^{\star},\ p^{\star})$
  };
  \node[colblock, seminovel, below=of cascade] (algo) {%
    \textbf{Algorithm 1A or 1B}\\[-1pt]
    \scriptsize calibrate $(w^{\dagger},\ b_{\epsilon}^{\star})$
  };
  \node[colblockio, semiruntime, below=of algo] (cache) {%
    \textbf{Calibration cache}\\[-1pt]
    \scriptsize key: scenario class; value: $(w^{\dagger},\ b_{\epsilon}^{\star})$
  };

  \node[phasehdr, below=8mm of cache] (runtimehdr) {Runtime loop (10 Hz)};
  \node[colblockio, semioffline, below=2mm of runtimehdr] (lookup) {%
    \textbf{Tick input + cache lookup}\\[-1pt]
    \scriptsize retrieve the calibrated pair for the active scenario class
  };
  \node[colblock, semianalysis, below=of lookup] (wssolve) {%
    \textbf{Weighted-sum solve}\\[-1pt]
    \scriptsize solve once at $w^{\dagger}$ to obtain $z_{\mathrm{ws}}^{\star}$
  };
  \node[coldecision, semidecision, below=7mm of wssolve] (check) {%
    Compliance preserved?\\[-1pt]
    $b_{\epsilon}(z_{\mathrm{ws}}^{\star}) = b_{\epsilon}^{\star}$
  };
  \node[hblockio, semiruntime, minimum width=3.0cm, text width=2.8cm,
        anchor=north] (accept) at ([yshift=-13mm, xshift=-18mm] check.south) {%
    \textbf{Accept}\\[-1pt]
    \scriptsize deploy trajectory
  };
  \node[hblock, semiwarn, minimum width=3.0cm, text width=2.8cm,
        anchor=north] (fallback) at ([yshift=-13mm, xshift=18mm] check.south) {%
    \textbf{Mismatch}\\[-1pt]
    \scriptsize log and rerun cascade
  };

  \draw[flow] (scenario.south) -- (cascade.north);
  \draw[flow] (cascade.south)  -- (algo.north);
  \draw[flow] (algo.south)     -- (cache.north);
  \draw[flow] (cache.south)    -- (lookup.north);
  \draw[flow] (lookup.south)   -- (wssolve.north);
  \draw[flow] (wssolve.south)  -- (check.north);
  \draw[flow] ([xshift=-2mm] check.south) |- node[arrowlbl, left=1pt, pos=0.25] {yes} (accept.north);
  \draw[feedback] ([xshift=2mm] check.south) |- node[arrowlbl, right=1pt, pos=0.25] {no} (fallback.north);

  \begin{scope}[on background layer]
    \node[groupbg, fit=(offlinehdr)(scenario)(cache)] {};
    \node[groupbg, fit=(runtimehdr)(lookup)(wssolve)(check)(accept)(fallback)] {};
  \end{scope}
\end{tikzpicture}

    \caption{Offline-online operational workflow of the LCP deployment. The offline phase calibrates and caches $(w^{\dagger}, b_{\epsilon}^{\star})$ per scenario class. The runtime loop performs one weighted-sum solve per tick, checks whether the resulting solution preserves the cached compliance pattern, and branches explicitly to either accepted execution or mismatch handling. Cf.~\cite[Fig.~B4]{lcp2025}.}
    \label{fig:state_machine}
\end{figure}

Fig.~\ref{fig:state_machine} places the framework into its operational context. The architectural contribution of~\cite{lcp2025} is the decomposition of the $L+1$-stage cascade into an offline calibration step, executed once per scenario class, and a single online weighted-sum solve per MPC tick. The decomposition is conditional on the runtime compliance check: at every tick, the binary compliance vector $b_{\epsilon}(z_{\mathrm{ws}}^{\star})$ at the operational solution is compared against the cached reference $b_{\epsilon}^{\star}$ from the offline lex cascade. A match certifies that the active-set structure has not drifted and the weighted-sum solution is admissible; a mismatch indicates structural drift and triggers either an online recalibration or a cascade fallback for the offending tick. The deployment described in Section~\ref{sec:deployment} realises both phases of this workflow, and the empirical match-rate evidence is reported in Section~\ref{sec:results}~(\S\ref{sec:results:theorem41}).
